% section: 発展パターン

\section{発展パターン}\label{節：発展パターン}
ここでは更に発展的なパターンについて紹介する.
扱われる漸化式は,一部の難関大において時々出題されるようなものである.
基本的に誘導がついた状態で出題されるため,ひと目見て解けるようになる必要はそこまで高くないが,知っていて損はない.
ここで追加する変数係数型は階比型と階差型の組合せ,
二重数列は階差と帰納法を二方向へ広げる見方である.
新しい型名を覚えるだけでなく,既習の操作がどこで働くかを確かめよう.
\subsection{変数係数の一次漸化式}\label{節：変数係数の一次漸化式}
  \sectionref*{節：階比型}では添字ごとの倍率を掛け合わせ,
  \sectionref*{節：階差等比複合型}では同じ更新をする量を引いた.
  ここでは,倍率と付加項がどちらも添字に依存する場合を,この二つの操作で解こう.

  \subsubsection{倍率をそろえて,変化を足す}
  まず
  \BookMathLayout{\[
    a_1=2,\qquad
    a_{n+1}=\frac{n+2}{n}a_n+(n+1)(n+2)
    \quad(n\ge1)
  \]}{\[
    \begin{aligned}
      a_1&=2,\\
      a_{n+1}&=\frac{n+2}{n}a_n+(n+1)(n+2)\\
      &\quad(n\ge1)
    \end{aligned}
  \]}
  を考える.係数の分子・分母に注目すると,
  \[
    \frac{(n+1)(n+2)}{n(n+1)}=\frac{n+2}{n}
  \]
  である.\textbf{隣り合う項の比が元の倍率になる量}として$h_n=n(n+1)$を選べる.
  これは\sectionref*{節：補助数列の設計}で目標から割る量を選んだ方法である.
  $b_n=\frac{a_n}{n(n+1)}$と置いて実際に計算すると,
  \[
    \frac{a_{n+1}}{(n+1)(n+2)}
      =\frac{a_n}{n(n+1)}+1
  \]
  となる.従って$b_1=1$, $b_{n+1}=b_n+1$より$b_n=n$であり,
  \[
    a_n=n^2(n+1)
  \]
  を得る.分母は$n\ge1$で非零である.初項と元の式へ代入しても確かめられる.

  \subsubsection{一般の手順と,割り算の条件}
  $a_{n+1}=p_na_n+q_n$ $(n\ge1)$に対して,
  $h_1=1$, $h_{n+1}=p_nh_n$とする.
  すべての$p_n$が非零なら,$h_n$も非零であり,
  \[
    h_n=\prod_{k=1}^{n-1}p_k
  \]
  である.空積は$1$とする.両辺を$h_{n+1}$で割ると,
  \[
    b_n=\frac{a_n}{h_n},\qquad
    b_{n+1}-b_n=\frac{q_n}{h_{n+1}}
  \]
  だから,階差を足して
  \[
    a_n=h_n\left(a_1+\sum_{k=1}^{n-1}\frac{q_k}{h_{k+1}}\right)
  \]
  を得る.先ほどは定数倍して$h_n=n(n+1)$を選んだが,
  その場合の$b_1$を正しく合わせれば同じ答えになる.

  \begin{bookpoint}
    \textbf{正規化の役割}\par
    倍率を取り除くと,残った付加項を足す問題になる.
    積・和の表示を得ることと,最後まで簡単な式で計算できることは別である.
  \end{bookpoint}
  もし$p_j=0$なら,$a_{j+1}=q_j$はそれ以前の項によらず決まる.
  零を含む$h_n$で割らず,零となる更新の前後で区間を分け,
  非零の倍率が続く範囲で同じ方法を使う.

  \subsubsection{同じ更新をする量を引く別解}
  先ほどの例では,$v_n=n^2(n+1)$が元の更新を満たすことを直接確かめられる.
  この候補は,正規化後の$b_n=n$から見つけたものである.
  $d_n=a_n-v_n$と置くと,
  \[
    d_{n+1}=\frac{n+2}{n}d_n
  \]
  となり,階比型へ戻る.初項を一般の$a_1=s$に変えると,
  \[
    a_n=n(n+1)\left(n+\frac{s-2}{2}\right)
  \]
  である.一つの解$v_n$と,同次の解の定数倍に分かれている.
  これは\sectionref*{節：非同次と解の組合せ}の見方そのものである.

  \begin{bookpoint}
    \textbf{二つの方法の選び方}\par
    隣り合う比が分かるなら割る因子を探す.
    同じ更新をする既知の量が見えるならそれを引く.
    どちらも,付加項と倍率を分けて扱うための操作である.
  \end{bookpoint}
  操作だけの比較は\sectionref*{節：操作を選ぶ練習},
  一般項まで求める練習は\bookproblemref{practice-connections}{1}{接続演習問1}へ進もう.

\subsection{ガウス型}\label{節：ガウス型}
  ここまでの漸化式では,差を取る,定数を引く,別の数列に置き換えるなどの操作から,既知の形へ帰着させてきた.
  ガウス記号を含む場合も,この方針は変わらない.
  ただし,変形の手掛かりとして,整数性,不等式,余り,偶奇,数え上げなどが必要になる.
  ガウス型では,整数問題で用いる考え方が前面に現れるのである.

  ガウス記号があるというだけで,一つの解法が決まるわけではない.
  本項では,\textbf{ガウス記号が式の中で何をしているか}に注目して,10種類の使い方を整理する.
  各解説の直後に代表例題を一つ置いた.
  同じ問題が複数の使い方に関係することもあるので,分類は排他的な区分ではなく,着眼点を選ぶための見取り図として使おう.

  \subsubsection{下準備と分類の見取り図}\label{節：ガウス型の下準備}
    ガウス記号$\lfloor x\rfloor$または$[x]$は,実数$x$を超えない最大の整数を表し,床関数と呼ばれる.
    整数$k$について,
    \[
      \lfloor x\rfloor=k\iff k\le x<k+1
    \]
    である.\,$x\ge0$では小数点以下を切り捨てた値に一致するが,
    負数では0に近づく切り捨てと異なる.例えば$\lfloor-1.5\rfloor=-2$である.
    天井関数$\lceil x\rceil$は$x$以上の最小の整数であり,
    \[
      \lceil x\rceil=k\iff k-1<x\le k
    \]
    で定義される.

    \BookFigFloorGraph

    式を読むときは,床関数が\textbf{項の値,添字,付加項のどれに作用するか}を最初に確かめよう.
    そのうえで,以下の使い方を考える.
    \begin{enumerate}[label=(\arabic*)]
      \item 整数性・不等式で整数部分を確定する.
      \item 同じ値を取る区間を群としてまとめる.
      \item 余りを取り出し,有限個の値と周期を調べる.
      \item 床関数の差で境界・条件・繰り上がりを検出する.
      \item 丸め誤差と係数の関係を使う.
      \item 整数への丸めで不動点と停止を調べる.
      \item 添字を二つに分割する.
      \item 商と余りで整数の桁を読む.
      \item 小数部分を残し,回転として読む.
      \item 条件を満たす整数や格子点の個数を数える.
    \end{enumerate}
    最初の6種類は主に値やその変化,次の2種類は添字,
    最後の2種類は小数部分や個数という意味に注目するものである.
    一般項を求めるだけでなく,周期や止まる値,何を数えているかを調べることにも意味がある.

    なお,整数$k$について$\lfloor x+k\rfloor=\lfloor x\rfloor+k$だが,
    実数同士では$\lfloor x+y\rfloor=\lfloor x\rfloor+\lfloor y\rfloor$とは限らない.
    数項を計算する実験は着眼点を探すために使い,予想は不等式や帰納法などで確かめよう.

    \subsubsection{整数性と不等式で整数部分を確定する}\label{節：ガウス型の整数部分の確定}
    \bookterm{floor}{ガウス記号}を見たら,まず中にある数が整数かどうかを考えよう.
    整数$k$と実数$x$に対して,
    \[
      \lfloor k+x\rfloor=k+\lfloor x\rfloor
    \]
    が成り立つ.例えば,整数$a_n$に対して$\lfloor a_n+\sqrt3\rfloor=a_n+1$である.
    ただし,これを使う前に,初項から整数が保たれることを確かめる必要がある.

    中身が整数との和でなくても,\,$k\le X<k+1$を示せば$\lfloor X\rfloor=k$と分かる.
    平方根なら,先に隣り合う平方数の間に挟むとよい.
    \textbf{\bookterm{floor-function}{床関数}の値を確定してから,既知の\bookterm{recurrence}{漸化式}として解く}のがこの使い方である.

    \noindent \bookblocklink{example-gauss-role-1}{problem}{answer}{【例題】}\par\nobreak
    \begin{enumerate}[label=(\arabic*),parsep=3pt,beginpenalty=10000,format=\bookitemlink{example-gauss-role-1}{problem}{answer}]
      \item \bookterm{sequence}{数列}を次のように定める.\bookterm{general}{一般項}を求めよ.
\BookMathLayout{            \[
              a_1=1,\qquad
              a_{n+1}=\left\lfloor\sqrt{a_n^2+4a_n+5}\right\rfloor
              \quad(n\ge1).
            \]}{            \[
              \begin{aligned}
              a_1&=1,\\
              a_{n+1}&=\left\lfloor\sqrt{a_n^2+4a_n+5}\right\rfloor
              \\              &\quad(n\ge1).
              \end{aligned}
            \]}
    \end{enumerate}

    \noindent \bookblocklink{example-gauss-role-1}{hint}{problem}{【方針】}\par\nobreak
    \begin{enumerate}[label=(\arabic*),parsep=3pt,beginpenalty=10000,format=\bookitemlink{example-gauss-role-1}{hint}{problem}]
      \item 整数$a_n$に対して,平方根の中身を$(a_n+2)^2$と$(a_n+3)^2$で挟みます.
            まず各項が非負整数であることを確かめましょう.
    \end{enumerate}

    \noindent \bookblocklink{example-gauss-role-1}{answer}{problem}{【解答】}\par\nobreak
    \begin{enumerate}[label=(\arabic*),parsep=3pt,beginpenalty=10000,format=\bookitemlink{example-gauss-role-1}{answer}{problem}]
      \item \[
              a_n=2n-1\quad(n\ge1).
            \]
    \end{enumerate}

    \noindent \bookblocklink{example-gauss-role-1}{solution}{problem}{【解法】}\par\nobreak
    \begin{enumerate}[label=(\arabic*),parsep=3pt,beginpenalty=10000,format=\bookitemlink{example-gauss-role-1}{solution}{problem}]
      \item 初項は非負整数であり,非負整数$a_n$から次の項も非負整数になる.
            よって帰納法により,すべての項が非負整数である.
            非負整数$x$について,
            \[
              (x+2)^2\le x^2+4x+5<(x+3)^2
            \]
            が成り立つ.左の差は1,右の差は$2x+4>0$である.
            平方根を取ると,
            \[
              x+2\le\sqrt{x^2+4x+5}<x+3
            \]
            なので,\bookterm{floor-function}{床関数}の値は$x+2$である.
            従って$a_{n+1}=a_n+2$となり,\sectionref{節：等差型}に帰着して$a_n=2n-1$を得る.
    \end{enumerate}

    \subsubsection{同じ値を取る区間を群としてまとめる}\label{節：ガウス型の群数列}
    \bookterm{floor-function}{床関数}は,中身が同じ二つの連続整数の間にある間,同じ値を取る.
    例えば,正整数$k,j$に対して,
    \[
      \lfloor\sqrt k\rfloor=j
      \iff j^2\le k\le(j+1)^2-1
    \]
    だから,値$j$が$2j+1$項続く.
    このようなまとまりを一つの群として数えるのが,群\bookterm{sequence}{数列}の考え方である.
    実際に$\lfloor\sqrt k\rfloor$は,
    \[
      \underbrace{1,1,1}_{3\text{個}},\,
      \underbrace{2,2,2,2,2}_{5\text{個}},\,
      \underbrace{3,3,3,3,3,3,3}_{7\text{個}},\,\ldots
    \]
    と並ぶ.

    \bookterm{difference}{階差}にこの形が現れたら,\sectionref{節：階差型}と同じように足す.
    \textbf{群の境界と,最後の群のどこまでを足すか}を区別しよう.
    似た形でも,和の上端が$n$か$n-1$かで境界が変わる.

    \noindent \bookblocklink{example-gauss-role-2}{problem}{answer}{【例題】}\par\nobreak
    \begin{enumerate}[label=(\arabic*),parsep=3pt,beginpenalty=10000,format=\bookitemlink{example-gauss-role-2}{problem}{answer}]
      \item \bookterm{sequence}{数列}を次のように定める.\bookterm{general}{一般項}を求めよ.
            \[
              a_1=2,\qquad
              a_{n+1}=a_n+\lfloor\sqrt{n+2}\rfloor
              \quad(n\ge1).
            \]
            必要ならば,\,$m=\lfloor\sqrt{n+1}\rfloor$を用いてよい.
    \end{enumerate}

    \noindent \bookblocklink{example-gauss-role-2}{hint}{problem}{【方針】}\par\nobreak
    \begin{enumerate}[label=(\arabic*),parsep=3pt,beginpenalty=10000,format=\bookitemlink{example-gauss-role-2}{hint}{problem}]
      \item \bookterm{difference}{階差}を足し,\,$\lfloor\sqrt k\rfloor$の和にします.
            \,$k=1,2$で値がともに1であることを使い,和の範囲をそろえましょう.
    \end{enumerate}

    \noindent \bookblocklink{example-gauss-role-2}{answer}{problem}{【解答】}\par\nobreak
    \begin{enumerate}[label=(\arabic*),parsep=3pt,beginpenalty=10000,format=\bookitemlink{example-gauss-role-2}{answer}{problem}]
      \item \,$m=\lfloor\sqrt{n+1}\rfloor$とおくと,
            \[
              a_n=m(n+2)-\frac{m(m+1)(2m+1)}6.
            \]
    \end{enumerate}

    \noindent \bookblocklink{example-gauss-role-2}{solution}{problem}{【解法】}\par\nobreak
    \begin{enumerate}[label=(\arabic*),parsep=3pt,beginpenalty=10000,format=\bookitemlink{example-gauss-role-2}{solution}{problem}]
      \item \bookterm{difference}{階差}を足すと,
            \[
              a_n=2+\sum_{k=3}^{n+1}\lfloor\sqrt k\rfloor
                  =\sum_{k=1}^{n+1}\lfloor\sqrt k\rfloor.
            \]
            \,$n=1$では最初の和を空の和として0とする.
            \,$N=n+1,\ m=\lfloor\sqrt N\rfloor$とおく.
            第$m-1$群までは全項,第$m$群は$k=m^2$から$N$までの$N-m^2+1$項を足すので,
            \begin{align*}
              a_n&=\sum_{j=1}^{m-1}j(2j+1)+m(N-m^2+1)\\
                 &=\frac{(m-1)m(4m+1)}6\\
                 &\quad+m(N-m^2+1)\\
                 &=m(N+1)-\frac{m(m+1)(2m+1)}6.
            \end{align*}
            \,$N=n+1$を戻せば\bookterm{general}{一般項}を得る.\,$m=1$でも最初の和は0となり,初項と一致する.

            \noindent \textbf{【別解】}\\
            \,$\lfloor\sqrt k\rfloor$は,\,$j^2\le k$を満たす正整数$j$の個数でもある.
            同じ組$(j,k)$を$j$ごとに数えると,
            \begin{align*}
              a_n&=\sum_{j=1}^{m}(N-j^2+1)\\
                 &=m(N+1)-\frac{m(m+1)(2m+1)}6.
            \end{align*}
            群ごとにまとめる方法と,何を数えているかを読み直す方法がつながっている.
    \end{enumerate}

    \subsubsection{余りを取り出し,値の範囲と周期を調べる}\label{節：ガウス型の余りと周期}
    整数$X$と正整数$m$に対して,
    \[
      X-m\left\lfloor\frac Xm\right\rfloor
    \]
    は$X$を$m$で割った余りであり,\,$0,1,\ldots,m-1$のいずれかである.
    これは\bookterm{floor-function}{床関数}の定義から,この式が0以上$m$未満の整数になるためである.
    大きな数を計算する式に見えても,\textbf{値が有限個に制限される}ことに注目しよう.

    \,$x\equiv y\pmod m$は,\,$x-y$が$m$の倍数であることを表す合同式である.
    次の値が$n$にも依存するなら,\,$a_n$だけでなく$n$の余りも一緒に追う.
    同じ組に戻り,その組から次の組が一意に決まれば,以降は繰り返す.
    有限個の状態を追う場合でも,初項から繰り返すとは限らず,途中から周期的になることもある.

    \noindent \bookblocklink{example-gauss-role-3}{problem}{answer}{【例題】}\par\nobreak
    \begin{enumerate}[label=(\arabic*),parsep=3pt,beginpenalty=10000,format=\bookitemlink{example-gauss-role-3}{problem}{answer}]
      \item \bookterm{sequence}{数列}を次のように定める.\bookterm{general}{一般項}と最小の正の周期を求めよ.
            \[
              a_1=1,\qquad
              a_{n+1}=X_n-4\left\lfloor\frac{X_n}{4}\right\rfloor,
            \]
            \[
              X_n=a_n^2+n^2+n\quad(n\ge1).
            \]
    \end{enumerate}

    \noindent \bookblocklink{example-gauss-role-3}{hint}{problem}{【方針】}\par\nobreak
    \begin{enumerate}[label=(\arabic*),parsep=3pt,beginpenalty=10000,format=\bookitemlink{example-gauss-role-3}{hint}{problem}]
      \item 次の項は$X_n$を4で割った余りです.
            \,$(a_n,n\bmod4)$の組を追うと,同じ状態への戻りを確かめられます.
    \end{enumerate}

    \noindent \bookblocklink{example-gauss-role-3}{answer}{problem}{【解答】}\par\nobreak
    \begin{enumerate}[label=(\arabic*),parsep=3pt,beginpenalty=10000,format=\bookitemlink{example-gauss-role-3}{answer}{problem}]
      \item \[
              a_n=
              \begin{cases}
                1 & (n\equiv0,1\pmod4),\\
                3 & (n\equiv2,3\pmod4).
              \end{cases}
            \]
            最小の正の周期は4である.
    \end{enumerate}

    \noindent \bookblocklink{example-gauss-role-3}{solution}{problem}{【解法】}\par\nobreak
    \begin{enumerate}[label=(\arabic*),parsep=3pt,beginpenalty=10000,format=\bookitemlink{example-gauss-role-3}{solution}{problem}]
      \item 初項は整数.整数$a_n$から$X_n$も次の項も整数になるので,各項は整数である.
            第2項以降は$0\le a_n<4$であり,
            \[
              a_{n+1}\equiv a_n^2+n(n+1)\pmod4
            \]
            となる.順に追うと,
            \[
              a_1,a_2,a_3,a_4,a_5=1,3,3,1,1.
            \]
            第1項と第5項では値が1,添字の余りが1で一致する.
            以降も同じ組から同じ次の組が決まるので,帰納法により$a_{n+4}=a_n$である.
            従って【解答】の\bookterm{general}{一般項}が得られる.
            第1項と第4項の値だけは同じでも,添字の余りが違うことに注意しよう.

            最小周期を確認する.1と2は$a_1\ne a_2,\ a_1\ne a_3$より周期ではなく,
            3も$a_2\ne a_5$より周期ではない.4は周期なので,最小の正の周期は4である.
    \end{enumerate}

    \subsubsection{床関数の差で境界・条件・繰り上がりを検出する}\label{節：ガウス型の境界と条件}
    正整数$n$と整数$m\ge2$に対して,
    \[
      \left\lfloor\frac nm\right\rfloor-
      \left\lfloor\frac{n-1}{m}\right\rfloor
      =\begin{cases}
         1 & (n\text{が}m\text{の倍数}),\\
         0 & (\text{それ以外})
       \end{cases}
    \]
    となる.商の整数部分が増える境界を検出しているのである.
    同様に,\,$\lfloor\sqrt n\rfloor-\lfloor\sqrt{n-1}\rfloor$は$n$が平方数のときだけ1になる.
    \textbf{式の形で条件を表し,条件に応じて変化する\bookterm{recurrence}{漸化式}を作れる}という使い方である.

    二つの数の和でも,整数の境界を越えたかを調べられる.
    小数部分を$\{x\}=x-\lfloor x\rfloor$と書けば,
    \begin{align*}
      &\lfloor x+y\rfloor-\lfloor x\rfloor-\lfloor y\rfloor\\
      &\qquad=\lfloor\{x\}+\{y\}\rfloor\in\{0,1\}.
    \end{align*}
    これは小数部分の和で繰り上がりが起きたときだけ1になる.
    \bookterm{floor-function}{床関数}は一般には和に分配できないが,その差には意味がある.

    \noindent \bookblocklink{example-gauss-role-4}{problem}{answer}{【例題】}\par\nobreak
    \begin{enumerate}[label=(\arabic*),parsep=3pt,beginpenalty=10000,format=\bookitemlink{example-gauss-role-4}{problem}{answer}]
      \item \bookterm{sequence}{数列}を次のように定める.\bookterm{general}{一般項}を求め,\bookterm{difference}{階差}の意味を説明せよ.
            \[
              a_0=0,\qquad a_n=a_{n-1}+d_n\quad(n\ge1),
            \]
            \[
              \begin{aligned}
                d_n={}&\left\lfloor\frac n4\right\rfloor
                       -\left\lfloor\frac{n-1}{4}\right\rfloor\\
                      &-\left\lfloor\frac n6\right\rfloor
                       +\left\lfloor\frac{n-1}{6}\right\rfloor.
              \end{aligned}
            \]
    \end{enumerate}

    \noindent \bookblocklink{example-gauss-role-4}{hint}{problem}{【方針】}\par\nobreak
    \begin{enumerate}[label=(\arabic*),parsep=3pt,beginpenalty=10000,format=\bookitemlink{example-gauss-role-4}{hint}{problem}]
      \item 各差を倍数かどうかの判定として読みます.
            和を取ると,途中の\bookterm{floor-function}{床関数}が打ち消し合います.
    \end{enumerate}

    \noindent \bookblocklink{example-gauss-role-4}{answer}{problem}{【解答】}\par\nobreak
    \begin{enumerate}[label=(\arabic*),parsep=3pt,beginpenalty=10000,format=\bookitemlink{example-gauss-role-4}{answer}{problem}]
      \item \[
              a_n=\left\lfloor\frac n4\right\rfloor
                  -\left\lfloor\frac n6\right\rfloor.
            \]
            \,$d_n$は4の倍数だけなら1,6の倍数だけなら$-1$,両方またはいずれでもないなら0である.
    \end{enumerate}

    \noindent \bookblocklink{example-gauss-role-4}{solution}{problem}{【解法】}\par\nobreak
    \begin{enumerate}[label=(\arabic*),parsep=3pt,beginpenalty=10000,format=\bookitemlink{example-gauss-role-4}{solution}{problem}]
      \item \bookterm{floor-function}{床関数}の差の性質から,\,$d_n$は【解答】に述べた条件を表す.
            また,\,$a_n=\sum_{k=1}^{n}d_k$であり,
            \begin{align*}
              \sum_{k=1}^{n}\left(
                \left\lfloor\frac k4\right\rfloor-
                \left\lfloor\frac{k-1}{4}\right\rfloor\right)
                &=\left\lfloor\frac n4\right\rfloor,\\
              \sum_{k=1}^{n}\left(
                \left\lfloor\frac k6\right\rfloor-
                \left\lfloor\frac{k-1}{6}\right\rfloor\right)
                &=\left\lfloor\frac n6\right\rfloor.
            \end{align*}
            両式の差から\bookterm{general}{一般項}を得る.\,$n=0$でも両辺は0である.
            条件を一つずつ場合分けして値を追う方法と,差を足して一度に求める方法がある.
    \end{enumerate}

    \subsubsection{丸め誤差と係数の関係からガウス記号を外す}\label{節：ガウス型の丸め誤差}
    \,$a_{n+1}=\lfloor\lambda a_n\rfloor$では,整数との和の場合と違い,\,$\lambda$の整数部分を取り出すだけでは解けない.
    代わりに,
    \[
      \varepsilon_n=\lambda a_n-a_{n+1},
      \qquad 0\le\varepsilon_n<1
    \]
    とおき,丸めで捨てた部分を残して計算しよう.
    \textbf{誤差そのものではなく,誤差の範囲から次の整数部分を確定する}のが狙いである.

    係数$\lambda$が2次方程式を満たす場合は,\,$\lambda^2$を$\lambda$と整数で表せることが手掛かりになる.
    ただし,最後に現れる誤差がどの整数の間にあるかを示す必要がある.
    どの無理数の係数でも同じ隣接3項間漸化式へ帰着できるわけではない.

    \noindent \bookblocklink{example-gauss-role-5}{problem}{answer}{【例題】}
    \begin{enumerate}[label=(\arabic*),parsep=3pt,beginpenalty=10000,format=\bookitemlink{example-gauss-role-5}{problem}{answer}]
      \item 数列を次のように定める.ガウス記号を含まない漸化式に帰着させ,一般項を求めよ.
            \[
              a_1=1,\qquad
              a_{n+1}=\lfloor(1+\sqrt2)a_n\rfloor
              \quad(n\ge1).
            \]
    \end{enumerate}

    \noindent \bookblocklink{example-gauss-role-5}{hint}{problem}{【方針】}
    \begin{enumerate}[label=(\arabic*),parsep=3pt,beginpenalty=10000,format=\bookitemlink{example-gauss-role-5}{hint}{problem}]
      \item \,$\lambda=1+\sqrt2$とおき,\,$\lambda^2=2\lambda+1$を使います.
            整数から正の1未満の誤差を引く形にして,床関数の値を確定しましょう.
    \end{enumerate}

    \noindent \bookblocklink{example-gauss-role-5}{answer}{problem}{【解答】}
    \begin{enumerate}[label=(\arabic*),parsep=3pt,beginpenalty=10000,format=\bookitemlink{example-gauss-role-5}{answer}{problem}]
      \item \[
              a_1=1,\quad a_2=2,\qquad
              a_{n+2}=2a_{n+1}+a_n-1.
            \]
            \,$\alpha=1+\sqrt2,\ \beta=1-\sqrt2$とおくと,
            \[
              a_n=\frac12+\frac{\alpha^n+\beta^n}{4}.
            \]
    \end{enumerate}

    \noindent \bookblocklink{example-gauss-role-5}{solution}{problem}{【解法】}
    \begin{enumerate}[label=(\arabic*),parsep=3pt,beginpenalty=10000,format=\bookitemlink{example-gauss-role-5}{solution}{problem}]
      \item \,$\lambda=1+\sqrt2>2$とおく.正整数$a_n$から次の項も正整数となる.
            また,\,$\lambda a_n$は無理数なので,
            \[
              \varepsilon_n=\lambda a_n-a_{n+1}
            \]
            は$0<\varepsilon_n<1$を満たす.
            \,$\lambda^2=2\lambda+1$を使うと,
            \begin{align*}
              \lambda a_{n+1}
                &=\lambda(\lambda a_n-\varepsilon_n)\\
                &=2a_{n+1}+a_n-(\lambda-2)\varepsilon_n.
            \end{align*}
            \,$0<\lambda-2=\sqrt2-1<1$なので,最後に引く数も0と1の間にある.
            \,$M_n=2a_{n+1}+a_n$は整数だから,
            \[
              M_n-1<\lambda a_{n+1}<M_n.
            \]
            従って$a_{n+2}=M_n-1$となり,【解答】の隣接3項間漸化式を得る.
            初期条件は$a_2=\lfloor1+\sqrt2\rfloor=2$である.

            定数項を消すため$b_n=a_n-\frac12$とおくと,
            \[
              b_{n+2}=2b_{n+1}+b_n,
              \qquad b_1=\frac12,\quad b_2=\frac32.
            \]
            \sectionref{節：隣接3項間漸化式型}で学んだ特性方程式は$x^2-2x-1=0$で,
            その解が$\alpha,\beta$である.
            \,$\alpha+\beta=2,\ \alpha^2+\beta^2=6$より,
            \[
              b_n=\frac{\alpha^n+\beta^n}{4}
            \]
            は初期条件と漸化式を満たす.
            二つの初期値から各項は一意に決まるので,この式が一般項である.
            丸め誤差を評価した後は,既習の線形の漸化式に戻ることができた.
    \end{enumerate}


    \subsubsection{整数への丸めで不動点と停止を調べる}\label{節：ガウス型の不動点と停止}
    丸めを繰り返す数列では,一般項だけでなく,どの整数で止まるかにも注目しよう.
    次の値が今の値と等しくなる値を不動点という.
    不動点に到達すると,その後は同じ値が続く.
    ただし,丸めを外した実数の漸化式と,不動点が同じとは限らない.

    整数$K$を引いて,
    \[
      b_{n+1}=\left\lfloor\frac{b_n}{2}\right\rfloor
    \]
    に帰着できる場合がある.
    正整数$m$に対して,
    \[
      \left\lfloor\frac{\lfloor x\rfloor}{m}\right\rfloor
      =\left\lfloor\frac xm\right\rfloor
    \]
    が成り立つので,繰り返すと分母をまとめられる.
    この性質を確かめよう.\,$k=\lfloor x\rfloor$を$m$で割り,\,$k=mq+r$と書くと,
    \,$0\le r\le m-1$であり,\,$x=mq+r+t\ (0\le t<1)$となる.
    これより$0\le r+t<m$だから,どちらの床関数も$q$になる.

    \textbf{初項の位置と,負の数の床関数の値}も重要である.
    例えば,\,$-1<x<0$なら$\lfloor x\rfloor=-1$であって0ではない.

    \noindent \bookblocklink{example-gauss-role-6}{problem}{answer}{【例題】}
    \begin{enumerate}[label=(\arabic*),parsep=3pt,beginpenalty=10000,format=\bookitemlink{example-gauss-role-6}{problem}{answer}]
      \item 非負整数$M$に対し,数列を次のように定める.
            \[
              a_1=M,\qquad
              a_{n+1}=\left\lfloor\frac{a_n+6}{2}\right\rfloor
              \quad(n\ge1).
            \]
            一般項と,最終的に一定になる値を求めよ.
    \end{enumerate}

    \noindent \bookblocklink{example-gauss-role-6}{hint}{problem}{【方針】}
    \begin{enumerate}[label=(\arabic*),parsep=3pt,beginpenalty=10000,format=\bookitemlink{example-gauss-role-6}{hint}{problem}]
      \item 6を引くと,床関数によって半分にする漸化式になります.
            初項が6より小さい場合は,負の数の床関数を使って考えましょう.
    \end{enumerate}

    \noindent \bookblocklink{example-gauss-role-6}{answer}{problem}{【解答】}
    \begin{enumerate}[label=(\arabic*),parsep=3pt,beginpenalty=10000,format=\bookitemlink{example-gauss-role-6}{answer}{problem}]
      \item \[
              a_n=6+\left\lfloor\frac{M-6}{2^{n-1}}\right\rfloor.
            \]
            有限回の操作の後で,\,$M<6$なら5,\,$M\ge6$なら6に固定される.
    \end{enumerate}

    \noindent \bookblocklink{example-gauss-role-6}{solution}{problem}{【解法】}
    \begin{enumerate}[label=(\arabic*),parsep=3pt,beginpenalty=10000,format=\bookitemlink{example-gauss-role-6}{solution}{problem}]
      \item 初項からすべての項が非負整数となる.
            \,$b_n=a_n-6$とおくと,整数6を床関数の外に出せるので,
            \begin{align*}
              b_{n+1}
                &=\left\lfloor\frac{b_n+12}{2}\right\rfloor-6\\
                &=\left\lfloor\frac{b_n}{2}\right\rfloor.
            \end{align*}
            床関数の合成の性質と帰納法により,
            \[
              b_n=\left\lfloor\frac{M-6}{2^{n-1}}\right\rfloor.
            \]
            \,$M\ge6$では,十分大きな$n$で中身が$[0,1)$に入り,床は0になる.
            \,$M<6$では,十分大きな$n$で中身が$(-1,0)$に入り,床は$-1$になる.
            これより【解答】の値に有限回で到達する.

            実際,整数$x$が不動点となる条件は,
            \[
              x\le\frac{x+6}{2}<x+1
              \iff 4<x\le6
            \]
            なので,不動点は5と6の二つである.
            初項によって到達する不動点が分かれることは,この後の\sectionref{節：初項依存型}ともつながる.
    \end{enumerate}


    \subsubsection{床関数で添字を二つに分割する}\label{節：ガウス型の添字の分割}
    ここからは,\bookterm{floor-function}{床関数}が\bookterm{sequence}{数列}の値ではなく,\textbf{どの項を参照するか}を指定する使い方を考える.
    正整数$n$に対して,
    \[
      \left\lfloor\frac{n+1}{2}\right\rfloor
      +\left\lfloor\frac n2\right\rfloor=n
    \]
    であり,二つの添字は$n$をできるだけ等しく分けた大きさになる.
    偶数では両方が同じ添字になり,奇数では隣り合う添字になることが手掛かりである.

    この形の問題では,まず偶奇で分けて\bookterm{floor-function}{床関数}を外し,それから\bookterm{difference}{階差}を取る方法がある.
    半分の添字との関係が得られれば,\,$1,\ 2\le n<4,\ 4\le n<8,\ldots$のように,
    2の累乗ごとの群で整理できる.
    定義の確認も必要で,例えば$n=1$で元と同じ添字を参照する式は,
    そのまま初項を決める式として使えないことがある.

    \noindent \bookblocklink{example-gauss-role-7}{problem}{answer}{【例題】}\par\nobreak
    \begin{enumerate}[label=(\arabic*),parsep=3pt,beginpenalty=10000,format=\bookitemlink{example-gauss-role-7}{problem}{answer}]
      \item \bookterm{sequence}{数列}を次のように定める.\bookterm{general}{一般項}を求めよ.
\BookMathLayout{            \[
              a_1=2,\qquad
              a_n=a_{\left\lfloor\frac{n+1}{2}\right\rfloor}
                  +a_{\left\lfloor\frac n2\right\rfloor}+n-1
              \quad(n\ge2).
            \]}{            \[
              \begin{aligned}
              a_1&=2,\\
              a_n&=a_{\left\lfloor\frac{n+1}{2}\right\rfloor}
                  +a_{\left\lfloor\frac n2\right\rfloor}+n-1
              \\              &\quad(n\ge2).
              \end{aligned}
            \]}
            必要ならば,\,$2^m\le n<2^{m+1}$を満たす整数$m\ge0$を用いてよい.
    \end{enumerate}

    \noindent \bookblocklink{example-gauss-role-7}{hint}{problem}{【方針】}\par\nobreak
    \begin{enumerate}[label=(\arabic*),parsep=3pt,beginpenalty=10000,format=\bookitemlink{example-gauss-role-7}{hint}{problem}]
      \item 偶奇で式を分け,\,$b_n=a_{n+1}-a_n$を考えます.
            \,$b_{2k}$と$b_{2k+1}$が$b_k+1$になることを示し,群ごとに足しましょう.
    \end{enumerate}

    \noindent \bookblocklink{example-gauss-role-7}{answer}{problem}{【解答】}\par\nobreak
    \begin{enumerate}[label=(\arabic*),parsep=3pt,beginpenalty=10000,format=\bookitemlink{example-gauss-role-7}{answer}{problem}]
      \item \,$2^m\le n<2^{m+1}$を満たす整数$m\ge0$に対して,
            \[
              a_n=(m+3)n-2^{m+1}+1.
            \]
    \end{enumerate}

    \noindent \bookblocklink{example-gauss-role-7}{solution}{problem}{【解法】}\par\nobreak
    \begin{enumerate}[label=(\arabic*),parsep=3pt,beginpenalty=10000,format=\bookitemlink{example-gauss-role-7}{solution}{problem}]
      \item \,$n\ge2$では二つの添字は1以上$n-1$以下なので,\bookterm{sequence}{数列}は順に定義できる.
            偶奇で分けると,\,$k\ge1$について,
            \[
              \begin{aligned}
                a_{2k}&=2a_k+2k-1,\\
                a_{2k+1}&=a_{k+1}+a_k+2k.
              \end{aligned}
            \]
            \,$b_n=a_{n+1}-a_n$とおくと,両式の差から$b_{2k}=b_k+1$.
            また,\,$a_{2k+2}=2a_{k+1}+2k+1$を使うと$b_{2k+1}=b_k+1$である.
            \,$a_2=5$より$b_1=3$なので,
            \[
              b_{2k}=b_{2k+1}=b_k+1,\qquad b_1=3.
            \]

            \,$2^m\le n<2^{m+1}$の範囲では$b_n=m+3$となる.
            \,$m=0$では$b_1=3$と一致する.
            次の群の添字を半分にして床を取ると前の群に入り,値が1増えるから,
            群番号$m$についての帰納法でこの式が示される.

            従って,完全な群と最後の群を分けて足すと,
            \begin{align*}
              a_n={}&2+\sum_{j=0}^{m-1}2^j(j+3)\\
                   &+(n-2^m)(m+3).
            \end{align*}
            ここで,
            \[
              \sum_{j=0}^{m-1}2^j=2^m-1,
            \]
            \[
              \sum_{j=0}^{m-1}j2^j=(m-2)2^m+2
            \]
            を使う.後者は,和を$T$とおき$2T-T$を計算すれば得られる.
            \,$m=0$でも空の和として両式は成り立つ.
            これを代入すると,
            \[
              a_n=(m+3)n-2^{m+1}+1
            \]
            を得る.添字を分割する問題から,\bookterm{difference}{階差}を通じて群\bookterm{sequence}{数列}へ戻ることができた.
    \end{enumerate}

    \subsubsection{商と余りで整数の桁を読む}\label{節：ガウス型の桁の抽出}
    正整数$b\ge2$に対し,非負整数$n$を$b$で割って,
    \[
      q=\left\lfloor\frac nb\right\rfloor,
      \qquad r=n-b\left\lfloor\frac nb\right\rfloor
    \]
    とおくと,\,$n=bq+r,\ 0\le r<b$である.
    十進法では,\,$q=\lfloor\frac n{10}\rfloor$が一の位を取り除いた整数,
    \,$r=n-10\lfloor\frac n{10}\rfloor$が一の位の数字になる.

    一般の$b$進法も,各桁の数字を0から$b-1$までとし,位を$b$倍ずつにする表し方である.
    このとき$q$は最後の桁を取り除く操作,$r$はその桁の数字に当たる.
    \textbf{添字を小さくする床関数と,桁を取り出す余りを組み合わせる}ことができる.
    前の添字分割は二つの項を足したが,ここでは一つの項へ戻りながら最後の桁を処理する.

    \noindent \bookblocklink{example-gauss-role-8}{problem}{answer}{【例題】}
    \begin{enumerate}[label=(\arabic*),parsep=3pt,beginpenalty=10000,format=\bookitemlink{example-gauss-role-8}{problem}{answer}]
      \item 数列を次のように定める.
            \[
              a_0=0,\qquad
              a_n=a_{\left\lfloor\frac n3\right\rfloor}
                  +n-3\left\lfloor\frac n3\right\rfloor
              \quad(n\ge1).
            \]
            \,$a_n$が$n$の3進表示とどのような関係を持つか説明し,\,$a_{14}$を求めよ.
    \end{enumerate}

    \noindent \bookblocklink{example-gauss-role-8}{hint}{problem}{【方針】}
    \begin{enumerate}[label=(\arabic*),parsep=3pt,beginpenalty=10000,format=\bookitemlink{example-gauss-role-8}{hint}{problem}]
      \item \,$n=3q+r\ (0\le r<3)$と書くと,漸化式は$a_n=a_q+r$になります.
            \,$q$の3進表示に最後の桁$r$を付けると$n$になることを使いましょう.
    \end{enumerate}

    \noindent \bookblocklink{example-gauss-role-8}{answer}{problem}{【解答】}
    \begin{enumerate}[label=(\arabic*),parsep=3pt,beginpenalty=10000,format=\bookitemlink{example-gauss-role-8}{answer}{problem}]
      \item \,$a_n$は$n$の3進表示の各桁の数字の和である.
            \,$14=(112)_3$より,
            \[
              a_{14}=1+1+2=4.
            \]
    \end{enumerate}

    \noindent \bookblocklink{example-gauss-role-8}{solution}{problem}{【解法】}
    \begin{enumerate}[label=(\arabic*),parsep=3pt,beginpenalty=10000,format=\bookitemlink{example-gauss-role-8}{solution}{problem}]
      \item \,$n\ge1$では$0\le q=\lfloor\frac n3\rfloor<n$なので,初期値$a_0$から各項を定義できる.
            \,$r=n-3q$とおけば$r$は0,1,2のいずれかで,
            \[
              a_n=a_q+r.
            \]
            \,$n$より小さな整数について桁の和との一致が分かっているとしよう.
            \,$n=3q+r$の3進表示は,\,$q$の表示を一桁ずらして最後に$r$を付けたものになる.
            その桁の和は$q$の桁の和に$r$を足した値であり,漸化式と一致する.
            \,$a_0=0$から,この考えを$n=1,2,\ldots$と順に適用すれば,すべての$n$について示される.
            \,$q=0$の場合も,桁の和を0として同じ議論が成り立つ.

            実際,\,$14=3\cdot4+2,\ 4=3\cdot1+1,\ 1=3\cdot0+1$なので,
            \[
              a_{14}=a_4+2=a_1+3=a_0+4=4.
            \]
            この問題では,一般項を長い式で書く前に,数列が何を表すかを理解することに意味がある.
    \end{enumerate}


    \subsubsection{小数部分を残し,区間内の回転として読む}\label{節：ガウス型の小数部分と回転}
    \,$\{x\}=x-\lfloor x\rfloor$は$x$の小数部分を表し,\,$0\le\{x\}<1$である.
    例えば,\,$\{-0.2\}=0.8$となる.
    ある量を加えてから整数部分を捨てる操作は,値を$[0,1)$へ戻す操作である.
    区間の両端をつないで円周と見れば,毎回同じ割合だけ進む回転として解釈できる.

    \textbf{小数部分を追うと,整数の境界を越えても運動の規則が残る}.
    途中で何回境界を越えたかは,\bookterm{floor-function}{床関数}で記録できる.
    有理数の割合なら繰り返しが生じるが,無理数の場合については,周期があるかを別に証明しよう.

    \noindent \bookblocklink{example-gauss-role-9}{problem}{answer}{【例題】}\par\nobreak
    \begin{enumerate}[label=(\arabic*),parsep=3pt,beginpenalty=10000,format=\bookitemlink{example-gauss-role-9}{problem}{answer}]
      \item \,$\alpha=\sqrt2-1$とし,\bookterm{sequence}{数列}を次のように定める.
            \[
              a_0=0,\qquad
              a_{n+1}=a_n+\alpha-\lfloor a_n+\alpha\rfloor
              \quad(n\ge0).
            \]
            \bookterm{general}{一般項}を求め,異なる添字の項が同じ値になるか調べよ.
    \end{enumerate}

    \noindent \bookblocklink{example-gauss-role-9}{hint}{problem}{【方針】}\par\nobreak
    \begin{enumerate}[label=(\arabic*),parsep=3pt,beginpenalty=10000,format=\bookitemlink{example-gauss-role-9}{hint}{problem}]
      \item 整数部分を捨てると小数部分だけが残ることから,\,$n\alpha$の小数部分を予想します.
            一致する二項があれば,添字の差と$\alpha$の積が整数になることを使いましょう.
    \end{enumerate}

    \noindent \bookblocklink{example-gauss-role-9}{answer}{problem}{【解答】}\par\nobreak
    \begin{enumerate}[label=(\arabic*),parsep=3pt,beginpenalty=10000,format=\bookitemlink{example-gauss-role-9}{answer}{problem}]
      \item \[
              a_n=n\alpha-\lfloor n\alpha\rfloor\quad(n\ge0).
            \]
            異なる添字の項は同じ値にならず,周期的ではない.
    \end{enumerate}

    \noindent \bookblocklink{example-gauss-role-9}{solution}{problem}{【解法】}\par\nobreak
    \begin{enumerate}[label=(\arabic*),parsep=3pt,beginpenalty=10000,format=\bookitemlink{example-gauss-role-9}{solution}{problem}]
      \item \,$n=0$では予想した式は$a_0=0$を与える.
            \,$a_n=n\alpha-\lfloor n\alpha\rfloor$とすると,整数との和の性質から,
            \begin{align*}
              \lfloor a_n+\alpha\rfloor
                &=\lfloor(n+1)\alpha\rfloor-\lfloor n\alpha\rfloor,\\
              a_{n+1}
                &=(n+1)\alpha-\lfloor(n+1)\alpha\rfloor.
            \end{align*}
            よって帰納法により\bookterm{general}{一般項}が示された.すべての項は$[0,1)$に属する.

            \,$a_{n+p}=a_n\ (p\ge1)$とすると,
            \[
              p\alpha=\lfloor(n+p)\alpha\rfloor-\lfloor n\alpha\rfloor
            \]
            が整数になり,\,$\alpha$が有理数となる.これは$\alpha=\sqrt2-1$が無理数であることに矛盾する.
            従って異なる添字の項は一致しない.

            同じ\bookterm{recurrence}{漸化式}で$\alpha$が有理数$\frac pq\ (0<p<q)$であり,\,$p,q$が互いに素なら,
            最小の正の周期は$q$となる.\,$q\alpha$が整数なので$q$項後に戻り,
            戻るまでの項数$T$は$T\alpha$が整数となる条件から$q$の倍数だからである.
            回転の見方と,有理数・無理数という整数問題の考え方が結びついている.
    \end{enumerate}

    \subsubsection{床関数を整数の個数として数える}\label{節：ガウス型の個数と格子点}
    \,$X\ge0$に対して$\lfloor X\rfloor$は,\,$1\le k\le X$を満たす整数$k$の個数である.
    例えば$\lfloor\frac nm\rfloor$は,\,$1,\ldots,n$に含まれる$m$の倍数の個数を表す.
    \bookterm{floor-function}{床関数}を含む和は,\textbf{何を数えているか}から意味を読み直せる.

    座標がともに整数である点を格子点という.
    例えば$\sum_{k=1}^{n}\lfloor ck\rfloor\ (c>0)$は,
    \,$x=k$を固定するたびに,\,$1\le y\le ck$を満たす整数$y$を数えている.
    同じ値の群で整理する方法とは別に,図形の中の点を数える方法としても理解できる.
    二つの使い方は重なり,前の群\bookterm{sequence}{数列}の別解も整数の組を数え直す方法であった.

    \noindent \bookblocklink{example-gauss-role-10}{problem}{answer}{【例題】}\par\nobreak
    \begin{enumerate}[label=(\arabic*),parsep=3pt,beginpenalty=10000,format=\bookitemlink{example-gauss-role-10}{problem}{answer}]
      \item \bookterm{sequence}{数列}を次のように定める.
            \[
              a_0=0,\qquad
              a_n=a_{n-1}+\left\lfloor\frac{5n}{2}\right\rfloor
              \quad(n\ge1).
            \]
            \bookterm{general}{一般項}を求め,\,$a_n$をある範囲の格子点の個数として説明せよ.
    \end{enumerate}

    \noindent \bookblocklink{example-gauss-role-10}{hint}{problem}{【方針】}\par\nobreak
    \begin{enumerate}[label=(\arabic*),parsep=3pt,beginpenalty=10000,format=\bookitemlink{example-gauss-role-10}{hint}{problem}]
      \item \,$\lfloor\frac{5k}{2}\rfloor=2k+\lfloor\frac k2\rfloor$と分けます.
            後者の和を偶奇で求め,\bookterm{floor-function}{床関数}が数えている整数$y$の範囲も確認しましょう.
    \end{enumerate}

    \noindent \bookblocklink{example-gauss-role-10}{answer}{problem}{【解答】}\par\nobreak
    \begin{enumerate}[label=(\arabic*),parsep=3pt,beginpenalty=10000,format=\bookitemlink{example-gauss-role-10}{answer}{problem}]
      \item \[
              a_n=n(n+1)+\left\lfloor\frac{n^2}{4}\right\rfloor.
            \]
            これは整数$x,y$で,
            \[
              1\le x\le n,\qquad 1\le y\le\frac{5x}{2}
            \]
            を満たす格子点の個数である.
    \end{enumerate}

    \noindent \bookblocklink{example-gauss-role-10}{solution}{problem}{【解法】}\par\nobreak
    \begin{enumerate}[label=(\arabic*),parsep=3pt,beginpenalty=10000,format=\bookitemlink{example-gauss-role-10}{solution}{problem}]
      \item \bookterm{difference}{階差}を足し,整数$2k$を\bookterm{floor-function}{床関数}の外に出すと,
            \begin{align*}
              a_n&=\sum_{k=1}^{n}\left\lfloor\frac{5k}{2}\right\rfloor\\
                 &=n(n+1)+\sum_{k=1}^{n}\left\lfloor\frac k2\right\rfloor.
            \end{align*}
            後ろの和は$n=2m$なら,
            \[
              0+1+1+2+2+\cdots+(m-1)+(m-1)+m=m^2,
            \]
            \,$n=2m+1$なら$m^2+m$である.
            従っていずれの場合も,
            \[
              \sum_{k=1}^{n}\left\lfloor\frac k2\right\rfloor
              =\left\lfloor\frac{n^2}{4}\right\rfloor
            \]
            となり,\bookterm{general}{一般項}を得る.\,$n=0$でも成立する.

            格子点としての意味も確かめよう.\,$x=k$と固定したとき,
            \,$1\le y\le\frac{5k}{2}$を満たす整数$y$は$\lfloor\frac{5k}{2}\rfloor$個である.
            \,$k=1,\ldots,n$について足した値が$a_n$だから,【解答】の範囲の点を数えている.
            \bookterm{sequence}{数列}の値,\bookterm{floor-function}{床関数}の和,図形の中の個数が同じ量を表している.
    \end{enumerate}


  \medskip
  \noindent \textbf{【着眼点の振り返り】}\\
    ガウス記号の値を不等式で確定する方法,群や余りで値を整理する方法,
    添字を分割する方法,小数部分や個数として読む方法を見てきた.
    そのどれを使うかは,記号の位置と,式の中で果たす役割から考えよう.
    余りと桁,群数列と数え上げのように,一つの問題から別の見方へつながることもある.

\subsection{隣接4項間漸化式型}\label{節：隣接4項間漸化式型}
隣接2項間や隣接3項間の漸化式を今まで扱った.
これをより広い範囲に適用してみよう.
例えば
\[
  a_{n+3}+pa_{n+2}+qa_{n+1}+ra_n=0
\]
のような漸化式を考える.
これは三階の同次線形漸化式なので,\sectionref*{節：一次独立・基底・次元}で述べたように,
初めの三項を自由に選ぶ解全体は3次元である.
ここでは一つの補助数列を先に求めることで,残りの計算を二階の式へ分ける.
\sectionref*{節：補助数列の設計}で述べたように,先に単純な更新を決め,それをする量を探そう.
3項をまとめた$d_n=a_{n+2}+\alpha a_{n+1}+\beta a_n$を作り,
$d_{n+1}=\gamma d_n$となることを目指す.
その条件は
\BookMathLayout{\[
  a_{n+3}+\alpha a_{n+2}+\beta a_{n+1}=\gamma \left(a_{n+2}+\alpha a_{n+1}+\beta a_{n}\right)
\]}{
\[
\begin{aligned}
a_{n+3}+\alpha a_{n+2}+\beta a_{n+1}\\
{}=\gamma\left(a_{n+2}+\alpha a_{n+1}+\beta a_n\right)
\end{aligned}
\]
}
である.両辺を左側に集め,元の漸化式と係数を比べると,
\[
  \alpha-\gamma=p,\qquad
  \beta-\gamma\alpha=q,\qquad
  -\gamma\beta=r
\]
となる.最初の二式から
\[
  \alpha=p+\gamma,\qquad
  \beta=q+p\gamma+\gamma^2
\]
なので,最後の式へ代入すると
\[
  \gamma^3+p\gamma^2+q\gamma+r=0
\]
を得る.この三次方程式の実数解を一つ選べば,$\alpha,\beta$も決まる.
三次方程式は少なくとも一つ実数解を持つが,計算しやすい値を見つけられるかは別の問題である.
以下の例題では因数分解や階差の形を手掛かりに選ぶ.
これによって,\,$d_n=a_{n+2}+\alpha a_{n+1}+\beta a_n$とおくと$d_{n+1}=\gamma d_n$が成り立つので,
\[
  a_{n+2}+\alpha a_{n+1}+\beta a_n=d_1\gamma^{\,n-1}
\]
という3項間漸化式が得られる.
ただし,\,\sectionref{節：等差二次特性方程式型}で扱ったのは右辺が定数$r$の場合であり,ここでは右辺が指数関数$d_1\gamma^{n-1}$になる点が異なる.
この形は\sectionref{節：階差等比複合型}の考え方を3項間に応用することで解くことができる.
ここでさらに同じ更新をする既知の量を引くと,右辺の指数関数を消せる場合がある.
例えば第1問では,$p=-3,q=3,r=-1$なので上の方程式は$(\gamma-1)^3=0$となり,
$\gamma=1,\alpha=-2,\beta=1$を得る.
補助数列が$d_n=a_{n+2}-2a_{n+1}+a_n$となる理由は,この係数の一致にある.
元の式が三階差の形であると気付けば,同じ量を階差からも作れる.

  次の数列の一般項を求めてみよう.

  \noindent \bookblocklink{example-four-term-1}{problem}{answer}{【例題】}
  \begin{enumerate}[label=(\arabic*),parsep=3pt,format=\bookitemlink{example-four-term-1}{problem}{answer}]
    \item $\displaystyle a_1=1,\,a_2=2,\,a_3=4$\\
          $a_{n+3}=3a_{n+2}-3a_{n+1}+a_n\quad(n\ge1)$
    \item $\displaystyle a_1=2,\,a_2=1,\,a_3=5$\\
          $a_{n+3}=2a_{n+2}+a_{n+1}-2a_n\quad(n\ge1)$
    \item $\displaystyle a_1=1,\,a_2=4,\,a_3=11$\\
          $a_{n+3}=5a_{n+2}-7a_{n+1}+3a_n\quad(n\ge1)$
  \end{enumerate}

  \noindent\bookblocklink{example-four-term-1}{hint}{problem}{【方針】}
  \begin{enumerate}[label=(\arabic*),parsep=3pt,format=\bookitemlink{example-four-term-1}{hint}{problem}]
    \item $d_n=a_{n+2}-2a_{n+1}+a_n$とおき,階差をさらにとると一定になることを利用します.
    \item $d_n=a_{n+2}-a_{n+1}-2a_n$とおくと$d_{n+1}=d_n$となります.初期条件から$d_n=0$を導き,3項間漸化式として解きます.
    \item $d_n=a_{n+2}-2a_{n+1}+a_n$とおくと$d_{n+1}=3d_n$となります.得られた右辺を満たす数列を引き,階差が一定の形に帰着させます.
  \end{enumerate}

  \noindent\bookblocklink{example-four-term-1}{answer}{problem}{【解答】}
  \begin{enumerate}[label=(\arabic*),parsep=3pt,format=\bookitemlink{example-four-term-1}{answer}{problem}]
    \item $\displaystyle a_n=\frac{n^2-n+2}{2}$
    \item $\displaystyle a_n=2^{n-1}+(-1)^{n-1}$
    \item $\displaystyle a_n=n-1+3^{n-1}$
  \end{enumerate}

  \noindent\bookblocklink{example-four-term-1}{solution}{problem}{【解法】}
  \begin{enumerate}[label=(\arabic*),parsep=3pt,format=\bookitemlink{example-four-term-1}{solution}{problem}]
    \item 漸化式を移項すると,
          \[
            a_{n+3}-3a_{n+2}+3a_{n+1}-a_n=0
          \]
          である.ここで
          \[
            d_n=a_{n+2}-2a_{n+1}+a_n
          \]
          とおくと,
          \[
            d_{n+1}-d_n=a_{n+3}-3a_{n+2}+3a_{n+1}-a_n=0
          \]
          より$d_n$は一定である.初期条件から$d_1=4-2\cdot2+1=1$なので,
          \[
            a_{n+2}-2a_{n+1}+a_n=1
          \]
          となる.さらに$e_n=a_{n+1}-a_n$とおくと$e_{n+1}-e_n=1$であり,$e_1=1$だから$e_n=n$である.したがって,
          \[
            a_n=a_1+\sum_{k=1}^{n-1}e_k
                =1+\sum_{k=1}^{n-1}k
                =\frac{n^2-n+2}{2}
          \]
          を得る.
    \item
          \[
            d_n=a_{n+2}-a_{n+1}-2a_n
          \]
          とおくと,
          \[
            d_{n+1}-d_n=a_{n+3}-2a_{n+2}-a_{n+1}+2a_n=0
          \]
          である.また,$d_1=5-1-2\cdot2=0$だから$d_n=0$であり,
          \[
            a_{n+2}-a_{n+1}-2a_n=0
          \]
          となる.特性方程式$x^2-x-2=0$の解は$2,-1$なので,
          \[
            a_n=A\,2^{n-1}+B(-1)^{n-1}
          \]
          とおける.初期条件より$A+B=2,\,2A-B=1$であるから$A=B=1$を得る.
    \item
          \[
            d_n=a_{n+2}-2a_{n+1}+a_n
          \]
          とおくと,与えられた漸化式から$d_{n+1}=3d_n$である.初期条件より$d_1=11-2\cdot4+1=4$なので,
          \[
            a_{n+2}-2a_{n+1}+a_n=4\cdot3^{n-1}
          \]
          となる.一方,$3^{n-1}$の2階差も$4\cdot3^{n-1}$である.そこで$b_n=a_n-3^{n-1}$とおくと,
          \[
            b_{n+2}-2b_{n+1}+b_n=0
          \]
          となる.$b_1=0,\,b_2=1$より$b_n=n-1$であるから,
          \[
            a_n=b_n+3^{n-1}=n-1+3^{n-1}
          \]
          を得る.
  \end{enumerate}

\subsection{虚数解型隣接3項間漸化式}\label{節：虚数解型隣接3項間漸化式}
  \BookFigComplexRoots

  今までの隣接3\,項間漸化式(\sectionref{節：隣接3項間漸化式型})では,特性方程式の解$\alpha,\beta$が異なれば$a_n=A\alpha^{n-1}+B\beta^{n-1}$という式自体は虚数解の場合でもそのまま成り立つ.
  ここでは,特性方程式が虚数解をもつ場合に,この式を三角関数を用いた実数の形に書き直す方法を考える.

  本項では,実数解で使った「特性根の冪を組み合わせる」流れを複素数へ広げる.
  \sectionref*{節：解を組み合わせる見方}で学んだ線形結合と初期条件の一意性を使い,
  実数の形への書き直しを確かめる.その後に,補助数列と加法定理から同じ形を見る.
  実数の係数$p,q$について,まず,
  \[
    a_{n+2}+pa_{n+1}+qa_n=0
  \]
  を考える.特性方程式は
  \[
    x^2+px+q=0
  \]
  である.

  これまで,特性方程式の解が異なる実数
  \[
    x=\alpha,\beta
  \]
  をもつ場合には,
  \[
    a_n=A\alpha^{n-1}+B\beta^{n-1}
  \]
  という形で考えた.本項の非実の根は零でないので,指数を一つずらした分を係数に吸収できる.
  以下では指数を$n$にした表し方も使う.

  では,特性方程式が実数解をもたず,
  \[
    x=\alpha,\overline{\alpha}
  \]
  という互いに共役な解をもつ場合には,一般項はどのような形になるだろうか.

  ここで「二つの冪で解全体を表せる」という理由を,前の共通説明につなげておこう.
  $p^2-4q<0$だから$q>0$であり,$\alpha,\overline\alpha$は異なる非零の根である.
  $u_n=\alpha^{n-1}$を元の式へ代入すると,
  \[
    u_{n+2}+pu_{n+1}+qu_n
    =\alpha^{n-1}(\alpha^2+p\alpha+q)=0
  \]
  となる.$v_n=\overline\alpha^{\,n-1}$も同じ式を満たす.
  それらの初めの二項は$(1,\alpha)$と$(1,\overline\alpha)$である.
  $Cu_n+Dv_n=0$なら$n=1,2$から
  \[
    C+D=0,\qquad C\alpha+D\overline\alpha=0
  \]
  を得る.$\alpha\ne\overline\alpha$なので$C=D=0$であり,二つの解は一次独立である.
  また,どんな初期条件もこの二つの係数で一意に合わせられる.
  従って,\sectionref*{節：一次独立・基底・次元}で述べた2次元の解全体の基底になっている.
  ここでは係数$C,D$も複素数として考えている.
  初めの二項を合わせれば,同じ更新からすべての項の一致が分かる.
  非実の根でも,実数解のときと同じ理由で冪の線形結合を使えるのである.

  例えば,
  \[
    x^2-2x+2=0
  \]
  を考えると,
  \[
    x=1\pm i
  \]
  である.したがって,漸化式
  \[
    a_{n+2}=2a_{n+1}-2a_n
  \]
  の特性方程式は虚数解をもつ.

  したがって,この例の一般項を
  \[
    a_n=A(1+i)^n+B(1-i)^n
  \]
  と書くことができる.しかし,\,$a_n$が実数列である場合,このままでは一般項の実数性が分かりにくい.そこで,複素数の極形式を利用する.

  一般に,特性方程式の解を
  \[
    \alpha=\rho(\cos\theta+i\sin\theta)=\rho e^{i\theta}
  \]
  とすると,共役解は
  \[
    \overline{\alpha}=\rho(\cos\theta-i\sin\theta)=\rho e^{-i\theta}
  \]
  である.したがって,
  \[
    \alpha^n=\rho^n e^{in\theta}
    =\rho^n(\cos n\theta+i\sin n\theta)
  \]
  および
  \[
    \overline{\alpha}^{\,n}
    =\rho^n e^{-in\theta}
    =\rho^n(\cos n\theta-i\sin n\theta)
  \]
  となる.

  ここで,$a_n$が実数列なら,共役を取っても値が変わらないので,
  \[
    A\alpha^n+B\overline\alpha^{\,n}
    =\overline B\alpha^n+\overline A\,\overline\alpha^{\,n}
  \]
  である.先ほど確かめた一次独立性から,係数の表し方は一意だから,
  $B=\overline A$となる.実数の解を選ぶ条件が,係数を共役にすることとして現れた.
  実際,\,$A=\dfrac{C-iD}{2},\ B=\dfrac{C+iD}{2}$(\,$C,D$は実数)とおくと,
\BookMathLayout{\begin{align*}
    A\alpha^n+B\overline{\alpha}^{\,n}
    &=\frac{C-iD}{2}\rho^n(\cos n\theta+i\sin n\theta)\\
    &\qquad+\frac{C+iD}{2}\rho^n(\cos n\theta-i\sin n\theta)\\
    &=\rho^n(C\cos n\theta+D\sin n\theta)
  \end{align*}}{
\[
\begin{aligned}
A\alpha^n+B\overline\alpha^{\,n}\\
{}=\frac{C-iD}{2}\rho^n(\cos n\theta+i\sin n\theta)\\
\quad+\frac{C+iD}{2}\rho^n(\cos n\theta-i\sin n\theta)\\
{}=\rho^n(C\cos n\theta+D\sin n\theta)
\end{aligned}
\]
}
  のように,虚部が打ち消し合って
  \[
    a_n=\rho^n(C\cos n\theta+D\sin n\theta)
  \]
  という実数の形に書き直せることが分かる.

  したがって,虚数解をもつ隣接3\,項間漸化式では,
  \[
    a_n=\rho^n(C\cos n\theta+D\sin n\theta)
  \]
  という形が基本的な一般項となる.ここで,\,$\rho$は特性根の絶対値,\,$\theta$はその偏角である.

  では,一般項
  \[
    a_n=\rho^n(C\cos n\theta+D\sin n\theta)
  \]
  に含まれる定数$C,D$をどのように決定すればよいだろうか.

  これまでの漸化式と同様に,初期条件を代入することで$C,D$を決定することができる.例えば,\,$a_1,a_2$が与えられているとする.一般項に$n=1,2$を代入すると,
  \[
    a_1=\rho(C\cos\theta+D\sin\theta)
  \]
  および
  \[
    a_2=\rho^2(C\cos2\theta+D\sin2\theta)
  \]
  を得る.

  これは$C,D$についての連立方程式である.したがって,この2つの方程式を解けば$C,D$が決定される.

  特に,$a_0,a_1$から始め,漸化式を$n\ge0$で使う場合には,より簡単になる.一般項に$n=0$を代入すると,
  \[
    a_0=C
  \]
  となる.したがって,
  \[
    C=a_0
  \]
  が直ちに得られる.

  さらに,\,$n=1$を代入すると,
  \[
    a_1=\rho(C\cos\theta+D\sin\theta)
  \]
  であるから,
  \[
    D=
    \cfrac{\frac{a_1}{\rho}-C\cos\theta}{\sin\theta}
  \]
  となる.ここで,\,$\sin\theta\neq0$であることに注意する.特性根が実数ではなく虚部をもつ場合には,\,$\theta$は$0$や$\pi$ではないため,
  \[
    \sin\theta\neq0
  \]
  が成り立つ.

  したがって,\,$a_0,a_1$が与えられている場合には,
  \[
    C=a_0,\qquad
    D=\cfrac{\frac{a_1}{\rho}-a_0\cos\theta}{\sin\theta}
  \]
  として$C,D$を求めることができる.

  例えば,先ほどの
  \[
    a_{n+2}=2a_{n+1}-2a_n
  \]
  を考える.特性根は
  \[
    1+i
    =\sqrt{2}
    \left(
      \cos\frac{\pi}{4}
      +i\sin\frac{\pi}{4}
    \right)
  \]
  であるから,
  \[
    \rho=\sqrt{2},\qquad
    \theta=\frac{\pi}{4}
  \]
  である.したがって一般項は
  \[
    a_n=(\sqrt{2})^n
    \left(
      C\cos\frac{n\pi}{4}
      +D\sin\frac{n\pi}{4}
    \right)
  \]
  となる.

  例えば初期条件
  \[
    a_0=1,\qquad a_1=2
  \]
  が与えられているとする.まず,
  \[
    C=a_0=1
  \]
  である.また,
  \[
    2=\sqrt{2}
    \left(
      C\cos\frac{\pi}{4}
      +D\sin\frac{\pi}{4}
    \right)
  \]
  であるから,
  \[
    2=\sqrt{2}
    \left(
      \frac{1}{\sqrt{2}}
      +\frac{D}{\sqrt{2}}
    \right)
    =1+D
  \]
  より,
  \[
    D=1
  \]
  を得る.

  よって,この場合の一般項は
  \[
    a_n=(\sqrt{2})^n
    \left(
      \cos\frac{n\pi}{4}
      +\sin\frac{n\pi}{4}
    \right)
  \]
  となる.

  このように,虚数解をもつ場合でも,基本的な流れは実数解の場合と変わらない.まず特性根から$\rho,\theta$を求め,
  \[
    a_n=\rho^n(C\cos n\theta+D\sin n\theta)
  \]
  とおき,初期条件を代入して$C,D$を決定すればよい.

  \noindent\textbf{【発想のつながり：倍率をそろえると加法定理が現れる】}\par
  特性根を極形式で書いたのは,一歩ごとの倍率と回転を読み取るためだった.
  解と係数の関係から$p=-2\rho\cos\theta,q=\rho^2$なので,元の式は
  \[
    a_{n+2}-2\rho\cos\theta\,a_{n+1}+\rho^2a_n=0
  \]
  である.今度はこの係数の形から補助数列を設計してみよう.
  初項$a_1$に合わせて指数を$n-1$にし,以下では別の係数$c,d$を使う.
  係数に$\rho$と$\rho^2$があるので,添字一つにつき$\rho$倍される量を基準にする.
  \sectionref*{節：補助数列の設計}の正規化を使い,
  \[
    b_n=\frac{a_n}{\rho^{n-1}}
  \]
  と置く.元の式を$\rho^{n+1}$で割ると
  \[
    b_{n+2}-2\cos\theta\,b_{n+1}+b_n=0
  \]
  となる.大きさの倍率を除いたことで,加法定理と同じ関係が残った.
  実際,
  \[
    \cos(x+\theta)+\cos(x-\theta)=2\cos\theta\cos x
  \]
  と,余弦を正弦に替えた同じ公式に$x=n\theta$を入れると,
  \[
    u_n=\cos((n-1)\theta),\qquad
    v_n=\sin((n-1)\theta)
  \]
  はどちらも$w_{n+2}-2\cos\theta\,w_{n+1}+w_n=0$を満たすと分かる.
  \sectionref*{節：線形結合}の重ね合わせの原理から,$cu_n+dv_n$も同じ関係を満たす.
  初めの二項は$(1,\cos\theta)$と$(0,\sin\theta)$であり,$\sin\theta\ne0$なので一次独立である.
  先ほどの複素数の冪に代わり,正弦・余弦が実数の解全体の基底になる.
  そこで
  \[
    b_n=c\cos((n-1)\theta)+d\sin((n-1)\theta)
  \]
  を候補にしよう.

  \noindent\textbf{【初期条件を合わせ,すべての項で確かめる】}\par
  候補の初めの二項は$c$と$c\cos\theta+d\sin\theta$である.
  元の初期条件から$b_1=a_1,b_2=\frac{a_2}{\rho}$なので,
  \[
    c=a_1,\qquad
    d=\cfrac{\frac{a_2}{\rho}-a_1\cos\theta}{\sin\theta}
  \]
  と選べば一致する.$\sin\theta\ne0$だから,この二つの係数は必ず決められる.
  候補も元の$b_n$も,直前の二項から同じ式で次の項を作る.
  従って初めの二項が一致すれば第三項も一致し,同じことを繰り返すとすべての項で一致する.
  初期条件と加法定理だけで,この候補が一般項であることを確かめられた.
  元へ戻すと,
  \BookMathLayout{\[
    a_n=\rho^{n-1}\{c\cos((n-1)\theta)+d\sin((n-1)\theta)\}
  \]}{\[
    \begin{aligned}
      a_n=\rho^{n-1}\bigl\{&c\cos((n-1)\theta)\\
                           &+d\sin((n-1)\theta)\bigr\}
    \end{aligned}
  \]}
  を得る.倍率をそろえ,加法定理に合う量を作り,初めの二項を合わせるのが解法の流れである.

  \noindent\textbf{【振動と周期性を区別する】}\par
  $C=D=0$なら零数列であり,以下では少なくとも一方が零でないとする.
  三角関数の部分には振動する形が現れるが,整数の添字で見た数列が周期的になるとは限らない.
  $\rho^{n}$は全体の大きさの倍率なので,
  \[
    \begin{cases}
      0<\rho<1 & \text{振幅が減衰する},\\
      \rho=1 & \text{振幅が一定である},\\
      \rho>1 & \text{振幅が増大する}
    \end{cases}
  \]
  と区別できる.
  $\rho=1$で$\frac{\theta}{2\pi}$が有理数なら,正整数$T$について$T\theta$を$2\pi$の整数倍にでき,
  加法定理から$a_{n+T}=a_n$となる.
  非零の解については,逆も成り立つ.
  複素数の冪の式で周期$T$を仮定すると,
  \[
    A(\alpha^T-1)\alpha^{n}
    +\overline A(\overline\alpha^{\,T}-1)\overline\alpha^{\,n}=0
  \]
  を得る.$n=1,2$の二つの式を使うと,$\alpha\ne\overline\alpha$だから両方の係数が零になる.
  非零の実数解では$A,\overline A$も零でないので,$\alpha^T=\overline\alpha^{\,T}=1$である.
  従って周期性には$\rho=1$と$\frac{\theta}{2\pi}\in\mathbb Q$の両方を確認する.
  三角関数があるというだけで周期性を結論しないようにしよう.

  次の数列の一般項を求めてみよう.

  \noindent \bookblocklink{example-complex-1}{problem}{answer}{【例題】}
  \begin{enumerate}[label=(\arabic*),parsep=3pt,format=\bookitemlink{example-complex-1}{problem}{answer}]
    \item $\displaystyle a_1=1,\,a_2=0,\quad a_{n+2}=-a_n\quad(n\ge1)$
\item \BookMathLayout{$\displaystyle a_1=0,\,a_2=1,\quad a_{n+2}=2a_{n+1}-2a_n\quad(n\ge1)$}{$\displaystyle a_1=0,\,a_2=1$\\ $\displaystyle a_{n+2}=2a_{n+1}-2a_n$\\ $\displaystyle(n\ge1)$}
    \item $\displaystyle a_1=1,\,a_2=2,\quad a_{n+2}=-4a_n\quad(n\ge1)$
\item \BookMathLayout{$\displaystyle a_1=1,\,a_2=0,\quad a_{n+2}=2a_{n+1}-4a_n\quad(n\ge1)$}{$\displaystyle a_1=1,\,a_2=0$\\ $\displaystyle a_{n+2}=2a_{n+1}-4a_n$\\ $\displaystyle(n\ge1)$}
  \end{enumerate}

  \noindent\bookblocklink{example-complex-1}{hint}{problem}{【方針】}
  \begin{enumerate}[label=(\arabic*),parsep=3pt,format=\bookitemlink{example-complex-1}{hint}{problem}]
    \item 初項から数項を計算して周期性を見つけ,一般項を予想します.その後,数学的帰納法で証明します.
    \item 特性根を極形式で表し,$\rho=\sqrt2,\,\theta=\dfrac{\pi}{4}$を用いて一般項を置きます.
    \item 特性根は$\pm2i$です.$\rho=2,\,\theta=\dfrac{\pi}{2}$として初期条件を代入します.
    \item 特性根を極形式で表し,$\rho=2,\,\theta=\dfrac{\pi}{3}$として係数を決定します.
  \end{enumerate}

  \noindent\bookblocklink{example-complex-1}{answer}{problem}{【解答】}
  \begin{enumerate}[label=(\arabic*),parsep=3pt,format=\bookitemlink{example-complex-1}{answer}{problem}]
    \item $\displaystyle a_n=\cos\frac{(n-1)\pi}{2}$
    \item $\displaystyle a_n=(\sqrt2)^{n-1}\sin\frac{(n-1)\pi}{4}$
    \item $\displaystyle a_n=2^{n-1}\left(\cos\frac{(n-1)\pi}{2}+\sin\frac{(n-1)\pi}{2}\right)$
\item \BookMathLayout{$\displaystyle a_n=2^{n-1}\left(\cos\frac{(n-1)\pi}{3}-\frac1{\sqrt3}\sin\frac{(n-1)\pi}{3}\right)$}{$\displaystyle a_n=2^{n-1}\Bigl(\cos\frac{(n-1)\pi}{3}$\\
$\displaystyle\hspace{3zw}-\frac1{\sqrt3}\sin\frac{(n-1)\pi}{3}\Bigr)$}
  \end{enumerate}

  \noindent\bookblocklink{example-complex-1}{solution}{problem}{【解法】}
  \begin{enumerate}[label=(\arabic*),parsep=3pt,format=\bookitemlink{example-complex-1}{solution}{problem}]
    \item 初項から数項を計算すると,
          \begin{align*}
            &a_1=1,\quad a_2=0,\quad a_3=-1,\quad a_4=0,\\
            &a_5=1,\quad a_6=0,\quad a_7=-1,\quad a_8=0
          \end{align*}
          となる.この並びから,
          \[
            a_n=\cos\frac{(n-1)\pi}{2}
          \]
          と予想できる.これを数学的帰納法で証明する.
          $n=1,2$では
          \[
            a_1=1=\cos0,\qquad a_2=0=\cos\frac{\pi}{2}
          \]
          となり,主張が成り立つ.
          ある$k\ge1$で,$n=k,k+1$に対して
          \[
            a_k=\cos\frac{(k-1)\pi}{2},\qquad
            a_{k+1}=\cos\frac{k\pi}{2}
          \]
          が成り立つと仮定する.漸化式と$-\cos x=\cos(x+\pi)$より,
          \[
            \begin{aligned}
              a_{k+2}
              &=-a_k\\
              &=-\cos\frac{(k-1)\pi}{2}\\
              &=\cos\left(\frac{(k-1)\pi}{2}+\pi\right)\\
              &=\cos\frac{(k+1)\pi}{2}
            \end{aligned}
          \]
          である.したがって,$n=k+1,k+2$でも主張が成り立つ.初めの2項で成り立つことと合わせ,数学的帰納法により
          \[
            a_n=\cos\frac{(n-1)\pi}{2}
          \]
          がすべての$n\ge1$で成り立つ.
    \item 特性方程式は
          \[
            x^2-2x+2=0
          \]
          であり,特性根は$1\pm i$である.したがって
          \[
            \rho=\sqrt2,\qquad\theta=\frac{\pi}{4}
          \]
          として,
\BookMathLayout{\[
            a_n=(\sqrt2)^{n-1}
            \left(
              C\cos\frac{(n-1)\pi}{4}
              +D\sin\frac{(n-1)\pi}{4}
            \right)
          \]}{
\[
\begin{aligned}
a_n&=(\sqrt2)^{n-1}\Bigl(C\cos\frac{(n-1)\pi}{4}\\
&\hspace{4zw}+D\sin\frac{(n-1)\pi}{4}\Bigr)
\end{aligned}
\]
}
          とおける.$a_1=0$より$C=0$であり,$a_2=1$より$D=1$である.よって,
          \[
            a_n=(\sqrt2)^{n-1}\sin\frac{(n-1)\pi}{4}
          \]
          を得る.
    \item 特性方程式は$x^2+4=0$であり,特性根は$2i,-2i$である.したがって
          \[
            a_n=2^{n-1}
            \left(
              C\cos\frac{(n-1)\pi}{2}
              +D\sin\frac{(n-1)\pi}{2}
            \right)
          \]
          とおける.$a_1=1$より$C=1$であり,$a_2=2$より$2D=2$,すなわち$D=1$である.よって,
          \[
            a_n=2^{n-1}
            \left(
              \cos\frac{(n-1)\pi}{2}
              +\sin\frac{(n-1)\pi}{2}
            \right)
          \]
          を得る.
    \item 特性方程式は
          \[
            x^2-2x+4=0
          \]
          であり,特性根は$1\pm\sqrt3\,i$である.したがって
          \[
            \rho=2,\qquad\theta=\frac{\pi}{3}
          \]
          として,
          \[
            a_n=2^{n-1}
            \left(
              C\cos\frac{(n-1)\pi}{3}
              +D\sin\frac{(n-1)\pi}{3}
            \right)
          \]
          とおける.$a_1=1$より$C=1$であり,$a_2=0$より
          \[
            0=2\left(\frac12+D\frac{\sqrt3}{2}\right)
          \]
          となるので$D=-\dfrac1{\sqrt3}$である.よって,
          \[
            a_n=2^{n-1}
            \left(
              \cos\frac{(n-1)\pi}{3}
              -\frac1{\sqrt3}\sin\frac{(n-1)\pi}{3}
            \right)
          \]
          を得る.
  \end{enumerate}

\subsection{非隣接2項間漸化式型}
  よく見られる非隣接二項間型の漸化式は
  \[
    a_{n+2}=f(a_{n})
  \]
  のような見た目のものである.
  このような場合には大きく分けて2つの解放がある.
  一つは$n=2k$のときと$n=2k-1$のときで場合分けして一般項を求める手法で,
  もう一方は0倍の$a_{n+1}$が存在すると考えて,隣接三項間漸化式に持ち込む方法である.

\subsection{二重数列と二方向の階差}\label{節：二重数列と二方向の階差}
  一つの番号に沿って項を並べる代わりに,行と列の二つの番号で値を並べてみよう.
  ここでは$n,m\ge0$を整数とし,$a_{n,m}$を数表の位置$(n,m)$の値とする.
  このように二つの添字をもつものを\textbf{二重数列}という.
  新しい万能の解法ではなく,階差と帰納法を二方向で使う見方を学ぶ.

  \subsubsection{距離を二乗し,二方向へ差を取る}
  原点から点$(n,m)$までの距離は
  \[ a_{n,m}=\sqrt{n^2+m^2} \]
  である.
  根号を取り除くため$b_{n,m}=a_{n,m}^2$とすると,
  一方向の差は$b_{n+1,m}-b_{n,m}=2n+1$となり,$m$によらない.
  従って,さらにもう一方向へ差を取れば零になる.

  この計算を広げ,境界の値
  \[
    b_{n,0}=n^2,\qquad b_{0,m}=m^2
  \]
  と,一定の実数$c$について
  \BookMathLayout{\[
    b_{n+1,m+1}-b_{n+1,m}-b_{n,m+1}+b_{n,m}=c
  \]}{\[
    \begin{aligned}
      b_{n+1,m+1}-b_{n+1,m}\\
      {}-b_{n,m+1}+b_{n,m}=c
    \end{aligned}
  \]}
  が与えられているとする.内側の値はこの式と境界から順に決まる.
  右辺を$b_{n+1,m+1}$について解けば,ほかの三つの位置は添字の和が小さい.
  $n+m$についての帰納法により一意性を確かめられる.

  横方向の差$d_{n,m}=b_{n+1,m}-b_{n,m}$を追うと,
  \[
    d_{n,m+1}-d_{n,m}=c,\qquad d_{n,0}=2n+1
  \]
  なので$d_{n,m}=2n+1+cm$である.これを横に足し戻せば,
  \[
    b_{n,m}=n^2+m^2+cnm
  \]
  となる.\textbf{二方向に一つずつ階差を取り,逆順に足し戻した}のである.
  $c$が一定という条件が,この短い計算を可能にしている.

  \subsubsection{図形と数え上げに戻る}
  二つの単位ベクトル$\boldsymbol u,\boldsymbol v$のなす角を$\theta$とすると,
  $c=2\cos\theta$の場合には
  \[
    b_{n,m}=|n\boldsymbol u+m\boldsymbol v|^2
  \]
  と読める.この図形的な解釈には$-2\le c\le2$が必要である.
  単に数表を作る式としてなら$c$は任意の実数でよいが,
  その値を距離の二乗と読むためには非負性も確認する.

  一方,格子点$(0,0)$から$(n,m)$へ,右と上へ一歩ずつ進む道の数を$A_{n,m}$とする.
  最後の一歩を分ければ,
  \[
    A_{n,m}=A_{n-1,m}+A_{n,m-1}\quad(n,m\ge1)
  \]
  であり,境界では$A_{n,0}=A_{0,m}=1$である.
  \BookFigLatticeRecurrence
  下の数表は,上と左の値を足して作ったものである.
  \[
    \begin{array}{c|rrrr}
      m\backslash n&0&1&2&3\\\hline
      0&1&1&1&1\\
      1&1&2&3&4\\
      2&1&3&6&10\\
      3&1&4&10&20
    \end{array}
  \]
  全$n+m$歩から上へ進む$m$歩の位置を選ぶので,
  \[
    A_{n,m}=\binom{n+m}{m}
  \]
  である.最後の一歩で分ける説明は,二項係数の加法公式の説明でもある.
  この公式と境界を使った帰納法による確認は,
  \bookproblemref{practice-connections}{3}{接続演習問3}で行おう.

  \begin{bookpoint}
    \textbf{添字を増やしても,確認することは同じ}\par
    どこが初期条件に当たる境界か,どの値から次が決まるかを先に確かめる.
    二方向の関係を別々に与える場合は,同じ位置を二通りに計算して矛盾しないことも必要である.
  \end{bookpoint}
  格子経路と二項係数の資料は参考文献のNIST DLMF §26.3へ.

